% !TeX root = ../main.tex
% =========================================================================
% 6. MARCO TEÓRICO REFERENCIAL
% Tres constructos deductivos (macro a micro), estándares internacionales
% y oraciones de transición al cierre de cada subsección.
% =========================================================================
\section{Marco teórico referencial}

Esta sección fundamenta conceptualmente la formulación del artefacto de software propuesto, estructurando los conocimientos de manera deductiva desde los paradigmas de interacción hasta los mecanismos de persistencia y los estándares de ingeniería de software que rigen su implementación. Los contenidos se organizan en torno a tres constructos esenciales: los agentes conversacionales en la atención comercial, los modelos generativos de lenguaje extenso, y las técnicas de recuperación aumentada sobre motores embebidos.

\subsection{Agentes conversacionales en la atención comercial}

Un agente conversacional es un sistema de software diseñado para emular la interacción humana mediante procesamiento de lenguaje natural, facilitando la ejecución de tareas o la provisión de información de manera interactiva \parencite{adamopoulou2020chatbots}. Históricamente, las primeras arquitecturas comerciales se fundamentaron en árboles de decisión y cotejo de patrones basados en reglas sintácticas fijas; si bien estos sistemas exhiben un comportamiento estrictamente predecible, sufren de una fragilidad intrínseca ante variaciones semánticas, errores tipográficos o expresiones idiomáticas no contempladas por los analistas \parencite{folstad2020users}. En el entorno del comercio electrónico conversacional, la percepción de naturalidad, comprensión contextual y empatía técnica condiciona directamente la satisfacción, la confianza del consumidor y su predisposición para culminar una transacción de compra \parencite{adam2021chatbots,ashfaq2020chatbot,sands2021who}. Asimismo, estudios empíricos en mercadotecnia digital demuestran que la pericia percibida del asistente y la transparencia en su condición de agente artificial inciden significativamente en la intención de compra y en la fidelización del cliente técnico \parencite{han2023conversational,luo2019frontiers}. Cuando se trata de suministros técnicos como materiales eléctricos, el usuario no solo busca conocer un precio, sino que requiere una consulta orientativa sobre compatibilidad de voltaje, calibres conductores y disponibilidad de marcas sustitutas. Esta complejidad lingüística y semántica hace inviable el uso exclusivo de árboles estáticos y motiva la integración de modelos generativos avanzados en el núcleo de procesamiento del agente.

\subsection{Modelos generativos de lenguaje extenso y arquitectura de agentes}

Los modelos de lenguaje extenso (LLM) contemporáneos se basan en la arquitectura Transformer, cuyo mecanismo de autoatención captura dependencias globales de significado a lo largo de secuencias textuales complejas \parencite{vaswani2017attention}. Entrenados sobre volúmenes masivos de datos textuales, estos modelos exhiben capacidades emergentes para la inferencia contextual, el seguimiento riguroso de directrices delimitadas en mensajes de sistema (\textit{system prompts}) y la adopción de roles comerciales especializados \parencite{wang2024survey}. No obstante, la naturaleza fundamentalmente probabilística de los LLM induce el fenómeno conocido como alucinación, en el cual el modelo produce aseveraciones sintácticamente coherentes pero fácticamente falsas o inexistentes en la realidad objetiva \parencite{ji2023hallucination,huang2023survey}. Recientes revisiones sistemáticas categorizan las alucinaciones en intrínsecas (contradicción directa de la fuente provista) y extrínsecas (invención de hechos no verificables), atribuyendo su origen a sesgos de preentrenamiento y fallas en la alineación paramétrica \parencite{tonmoy2024comprehensive}. En un entorno industrial o comercial de distribución eléctrica, la invención de un precio unitario, una capacidad de interrupción inexistente o una certificación técnica errónea acarrea graves perjuicios comerciales y legales. Por consiguiente, resulta imperativo gobernar el comportamiento generativo mediante verificación de cadenas de conocimiento y mecanismos de anclaje determinista (\textit{grounding}) \parencite{dhuliawala2023chain}, conocidos formalmente como generación aumentada por recuperación.

\subsection{Recuperación aumentada, bases de datos embebidas y estándares de ingeniería}

El paradigma de generación aumentada por recuperación (\textit{Retrieval-Augmented Generation}, RAG) constituye una metodología híbrida que desacopla la memoria paramétrica del modelo de lenguaje de la fuente fáctica de verdad, consultando un almacén de datos estructurado o no estructurado antes de sintetizar la respuesta final \parencite{guu2020realm,lewis2020retrieval,borgeaud2022improving,gao2023retrieval}. Si bien propuestas recientes exploran mecanismos avanzados de auto-reflexión y re-evaluación iterativa (\textit{Self-RAG}) \parencite{asai2024selfrag}, tales enfoques imponen una penalización severa de latencia y costo computacional. En el contexto de catálogos de materiales eléctricos, donde las entidades poseen atributos inequívocos como código de parte, amperaje, marca y precio de lista, la recuperación basada en esquemas relacionales e índices lexicográficos presenta mayor precisión y predictibilidad que las representaciones vectoriales densas, evitando discrepancias de redondeo y falsos positivos semánticos. 

Para satisfacer la restricción de hardware en entornos de bajo presupuesto computacional, el empleo de SQLite como motor embebido ofrece transacciones ACID completas dentro del mismo espacio de proceso de la aplicación, suprimiendo la sobrecarga de comunicación interproceso por sockets de red y el consumo de memoria asociado a servicios de bases de datos independientes \parencite{gaffney2022sqlite}. Al habilitar el modo de registro de escritura anticipada (\textit{Write-Ahead Logging}, WAL), el motor soporta lecturas simultáneas altamente concurrentes sin bloquear las escrituras de actualización de catálogo o registro de sesiones. Asimismo, la estructuración del backend bajo principios de microservicios independientes y patrones \textit{cloud-native} garantiza un despliegue desacoplado, mantenible y escalable con huella mínima de memoria en servidores virtuales de entrada \parencite{kratzke2017understanding,newman2021building,taibi2020patterns}.

Asimismo, la construcción del artefacto se gobierna bajo el estándar internacional de calidad de producto software \textbf{ISO/IEC 25010:2011}, priorizando las características de adecuación funcional, eficiencia de desempeño (tiempo de respuesta y utilización de recursos), fiabilidad y seguridad \parencite{iso25010}. La seguridad de la información en el canal de mensajería se garantiza mediante el protocolo \textbf{TLS 1.3} (RFC 8446) para el cifrado en tránsito \parencite{rescorla2018tls}, complementado con la verificación de firmas criptográficas HMAC-SHA256 exigida por WhatsApp Cloud API para autenticar la legitimidad de cada webhook entrante. Estos fundamentos teóricos, metodológicos y normativos proporcionan el marco de referencia indispensable sobre el cual se articula la propuesta de desarrollo técnico descrita en la siguiente sección.
